\section{Audited low-radius extensions after the \texorpdfstring{$e/2$}{e/2} frontier}
\label{sec:p12-low-radius-post-r23}

Put
\[
\varepsilon_{\max}:=c-T=\frac12\log\frac54,
\qquad
\rho:=\varepsilon_{\max}-\delta
=\frac12\log\frac{10}{9}.
\]
In the mixed strip write as before
\[
\sigma:=S-T,
\qquad
\varepsilon:=T_0-T,
\qquad
0<\sigma<\varepsilon<\varepsilon_{\max}.
\]

\begin{theorem}[A15.1b2e: global \(\rho\)-descent]
\label{thm:p12-rho-descent}
Let \(2a<T_0<c\). For every
\[
\rho\le R<T<S<T_0
\]
one has
\[
\boxed{
\ker L_{R,S,T_0}^{\{a,b,2a\}}=\{0\}.
}
\]
\end{theorem}

\begin{proof}
The final audited reassembly uses the exact partition
\[
\begin{array}{ll}
\text{A:}&R\ge e/2,\\
\text{B:}&R<e/2,\ \sigma\le R,\\
\text{C:}&R<e/2,\ R<\sigma<e/2,\\
\text{D:}&R<e/2,\ R<\sigma,\ \sigma\ge e/2.
\end{array}
\]
The partition is disjoint and exhaustive, including the boundary
assignments \(R=e/2\) to A, \(\sigma=R\) to B, and \(\sigma=e/2\) to D.

Case A is Theorem~\ref{thm:p12-b2d-full} together with the already
established higher-radius strata. Case B is the independently audited
restricted-tail closure of Round~19. Case D is the independently audited
Region-D closure of Round~20.

In Case C, Round~18 proves
\[
h=0\quad\text{a.e. on }(R,\sigma).
\]
Together with the original support cutoff this gives
\[
h=0\quad\text{a.e. on }(0,\sigma).
\]
The independently audited Round~17 full-tail lemma then gives
\[
H(t)=l(t)=0\quad\text{for a.e. }0<t<\sigma.
\]
Hence the same kernel vector is supported in \((\sigma,T)\). Applying
Theorem~\ref{thm:p12-b1} with the legal effective radius
\(R_{\rm eff}=\sigma\) and \(S_{\rm eff}=T\) yields \(h=0\).

The complete four-way composition was independently reviewed after
Round~20 and returned GREEN. This proves the theorem.
\end{proof}

\begin{theorem}[A15.1b2f: all-radius restricted-tail injectivity]
\label{thm:p12-restricted-tail-all-r}
Let \(2a<T_0<c\), \(0<R<T\), and
\[
0<\sigma\le R,
\qquad
\sigma<\varepsilon<\varepsilon_{\max}.
\]
Then
\[
\boxed{
\ker L_{R,T+\sigma,T_0}^{\{a,b,2a\}}=\{0\}.
}
\]
\end{theorem}

\begin{proof}
For \(R\ge e/2\) this follows from the established high-radius mixed-strip
result. Assume therefore \(0<R<e/2\).

For \(x\in(R,a)\), the condition \(x>R\ge\sigma\) puts the mixed-tail
term at \(T+x\) above \(S=T+\sigma\). Hence the lower E-equation is the
same endpoint equation as in the repaired A14.3a argument, and the lower
circle gives
\[
h=0\quad\text{a.e. on }(R,a).
\tag{R22.1}
\]
For \(0<t<\varepsilon\), the horizon source \(u=T+t\) gives
\[
p h(a+t)+r h(e+t)+q h(t)=0.
\]
Because \(\varepsilon<e\) and \(e+\varepsilon<2e<a\), both \(h(t)\)
and \(h(e+t)\) vanish by support or by (R22.1). Thus
\[
h(a+t)=0\quad(0<t<\varepsilon).
\]
Equivalently, with \(l(z)=h(T-z)\),
\[
l(z)=0\quad(a-\varepsilon<z<a).
\]
Since
\[
a-\varepsilon_{\max}>d,
\qquad
d-\varepsilon_{\max}>e/2>R\ge\sigma,
\]
the high-reflection identity on this seed strip is free of mixed-tail
contamination and transports the zero strip to
\[
l(w)=0\quad(0<w<\varepsilon).
\]
Now use the unconditional paired-source identity P1,
\[
H(t)+l(t)+\frac{2r}{p}H(d-t)=0.
\]
For \(0<t<\sigma\),
\[
d-t>d-\sigma\ge d-R>d-e/2>e/2>\sigma,
\]
so \(H(d-t)=0\), while \(l(t)=0\). Hence \(H(t)=0\) on the whole tail
\((0,\sigma)\). The same kernel vector is therefore supported in
\((R,T)\), and Theorem~\ref{thm:p12-b1} gives \(h=0\).

The proof uses no lower bound \(R\ge\rho\). Its retained structural
verifier and the four constant comparisons were independently rerun after
Round~22 and returned GREEN.
\end{proof}

\begin{proposition}[A15.1b2g: exact local overlap seed below \(\rho\)]
\label{prop:p12-r23-c23}
Define
\[
\omega:=\frac e2-\rho=\frac14\log\frac{27}{25},
\qquad
\eta:=e-2\delta,
\qquad
\chi:=3\delta-e,
\qquad
\kappa:=e-\delta.
\]
Suppose
\[
0<R<x<\sigma<\varepsilon<\varepsilon_{\max}
\]
and
\[
\boxed{
\begin{aligned}
&\omega<R,\\
&\eta<x<\chi,\\
&R+x<\delta<\sigma+x<\kappa,\\
&\sigma-x<\eta,\\
&\kappa<\varepsilon+x,\\
&x+\eta<\varepsilon.
\end{aligned}}
\tag{C23}
\]
Then every kernel vector satisfies
\[
\boxed{h(x)=0}
\]
for a.e. \(x\) in this cell.
\end{proposition}

\begin{proof}
The Round~23 certificate uses 42 horizon-legal affine sources
\(u=sx+me+n\delta\), split into 21 sources on each of the sheets
\(s=-1\) and \(s=+1\). Reconstructing every row directly from
\[
Lh(u)=p[h(u-a)-h(u+a)]
+r[h(u-b)-h(u+b)]
+q[h(u-T)-h(u+T)]
\]
with only odd reflection and the support cutoff gives exactly 42 live
visibility coordinates throughout the open cell (C23). Hence one obtains a
square matrix \(M_{42}\).

Exact symbolic elimination gives
\[
\det M_{42}=-p^{14}r^4F_-F_+.
\]
With
\[
\beta=q/p,
\qquad
v=(r/p)^2,
\]
the two degree-12 factors have the normalized form
\[
F_\pm=p^{12}(A\pm C),
\qquad
C=2\beta^3v^2(\beta^2-1)(2\beta^2+v-2).
\]
For the actual P12 weights
\[
\beta=2^{-3/4},
\qquad
v=\frac{\log3}{\log2}\sqrt{\frac8{27}},
\]
an exact rational interval computation gives
\[
A-C\in(-0.20661356870806363,-0.2066135687080636),
\]
\[
A+C\in(-0.0325699133677542,-0.032569913367754194).
\]
Thus both factors are strictly nonzero and
\[
\det M_{42}\ne0.
\]
Therefore every live coordinate vanishes, in particular \(h(x)=0\).
The raw reconstruction, determinant factorization, normalization, and a
separately written Fraction-based interval proof were independently rerun
after the repaired Round~23 commit and returned GREEN.
\end{proof}

\begin{remark}[Scope of the Round-23 seed]
Proposition~\ref{prop:p12-r23-c23} is local. It does not imply a global
statement for \(R\ge\omega\), and it does not close the whole overlap region
below \(\rho\).
\end{remark}

\begin{remark}[Audit provenance]
The detailed retained audit records are Round~21
(`8dbc594317f9ea993ea2f4b4256d61b3e71cc975`), Round~22 promotion
(`9ffa7aaf9903f411a589eb0fbc590f65f1eb9f60`), and Round~23 promotion
(`aa189c972961a90e364cb7b1a8ad64f02ab09bfd`). The Round~23 text payload was
repaired transparently in commit
`cddc6c5793f18a3dc02f65697063b395c9f7d897`; no history was rewritten.
\end{remark}
